\section{General maxmin representation: complete proofs}
\label{sec:annex-p17-maxmin-representation}

\begin{annexitem}{annex:p17.mmeu.utility}{thm:p17-mmeu-utility}{Mixture continuity, finite calibration and the arbitrary-outcome patch}
\begin{proof}
A.1 implies that \(f\succeq g\) is equivalent to the
failure of \(g\succ f\). Applying A.2 in both orientations
therefore preserves weak comparisons and indifference as
well as strict comparisons. On constants, its common
act may be any lottery, so this is full lottery
independence. Finite indifferent replacements and the
strict ordering of the segment
\(m(t)=tH+(1-t)L\), \(H\succ L\), follow from
\cref{thm:p17-mixture-comparison}, whose independence-only
proof applies on the finite union of the relevant supports.

We supply the calibration using the present A.3.
If \(H\succeq y\succeq L\), let
\(q=\sup\{t\in[0,1]:m(t)\preceq y\}\).
This is the supremum of a nonempty downward set.
The endpoint cases \(y\sim L,H\) give \(q=0,1\).
Otherwise \(H\succ y\succ L\), and A.3 gives an
interior mixture below \(y\) and an interior mixture
above \(y\), so \(0<q<1\).
If \(m(q)\succ y\), A.3 applied to
\(m(q)\succ y\succ L\) gives some
\(m(t)\succ y\) with \(t<q\), contradicting the
downward-set definition of the supremum.
If \(y\succ m(q)\), apply A.3 to
\(H\succ y\succ m(q)\); it gives an \(m(t)\prec y\)
with \(t>q\), another contradiction.
Thus \(y\sim m(q)\), uniquely because the segment
is strictly ordered.

On a finite prize set, choose best and worst certain
prizes. If all certain prizes are indifferent, finite
mixture replacement makes all its lotteries indifferent.
Otherwise calibrate each prize \(x\) as above, obtaining
a number \(v(x)\in[0,1]\).
Successively replacing the finitely many prizes in a
lottery by their indifferent binary mixtures gives
\[
             y\sim m\!\left(\sum_x y(x)v(x)\right).
\]
Thus this finite expectation represents all its lottery
comparisons. A second affine representation must give
the same calibrated numbers after normalizing the best
and worst prizes, proving finite affine uniqueness.
This derivation used A.3, not neighborhood continuity
on an unspecified outcome topology.

If the global constant-lottery order is nontrivial,
there are two strictly ranked certain prizes \(H,L\):
otherwise every finite-support lottery would be
indifferent by the preceding replacement argument.
For each finite prize set containing \(H,L\), construct
its finite representation and normalize it by
\(v(H)=1,v(L)=0\), even if these are not that set's
extreme prizes. On nested finite sets the
representations agree, by finite affine uniqueness.
For any \(x\in X\), define \(v(x)\) using
\(\{H,L,x\}\). For any finite collection of prizes,
adjoining \(H,L\) shows that these assignments agree
with its finite representation. Therefore
\(u(y)=\sum_x y(x)v(x)\) represents the preference
on all of \(Y\) and is affine. The same patch proves
global uniqueness up to one positive affine change.
No boundedness of \(v\) follows or is required.

If all constants are indifferent, A.4 in both
directions makes every pair of acts indifferent.
Thus A.6 implies the nontrivial case above.
Replacing its normalization by \(4u-2\) gives
values \(2,-2\) at \(H,L\); affine mixtures imply
\([-2,2]\subseteq K\), giving the stated interior.
In the trivial case a constant utility represents
the order, and any other affine representation of
that constant order is constant, hence related by
a positive affine change.
\end{proof}
\end{annexitem}

\begin{annexitem}{annex:p17.mmeu.certainty}{thm:p17-mmeu-certainty}{Every finite-step certainty equivalent}
\begin{proof}
The finitely many lottery values of \(f\) have a
best and a worst value under the complete order.
A.4 puts \(f\) between the corresponding constant
acts. In the nondegenerate case, include the two
fixed strictly ranked reference lotteries when
selecting bounds, ensuring strictly ordered bounds
without losing either inequality.
The calibration proof in the preceding item applies
with \(f\) as the middle object: its two A.3
applications mix only the constant upper and lower
objects, while the strict ordering of the constant
segment is already known. It gives \(f\sim y^*\).
If all constants are indifferent, A.4 already
gives this conclusion for every \(y\).

Any two such \(y\)'s are indifferent by
transitivity, hence have equal utility.
Define \(J(f)\) to be that common number.
Replacing \(f,g\) by their indifferent constants
shows \(f\succeq g\) exactly when \(J(f)\geq J(g)\).
It agrees with \(u\) on constant acts.
Any numerical functional representing this same
order and agreeing on constants must assign
that number to \(f\), proving uniqueness.
Strict segment endpoints give a unique coefficient;
if the two endpoints are indifferent, every
coefficient has the same utility, explaining
the distinction stated in the result.
\end{proof}
\end{annexitem}

\begin{annexitem}{annex:p17.mmeu.simple.functional}{thm:p17-mmeu-simple-functional}{All steps in Lemma 3.3}
\begin{proof}
For a simple function \(a\) with values in \(K\),
choose a lottery of the corresponding utility
on each cell of a finite partition representing
\(a\), and call the resulting act \(f_a\).
If two choices give the same utility at every
state, A.4 in both directions makes them
indifferent. Thus \(I_0(a)=J(f_a)\) is
well defined and monotone.
For \(c\in K\), \(I_0(c\mathbf1)=c\).
The certainty equivalent of \(f_a\) and A.2
give, for \(c\in K\) and \(0<t<1\),
\[
       I_0(ta+(1-t)c\mathbf1)=tI_0(a)+(1-t)c.
\]
The endpoints hold directly. Taking \(c=0\)
proves \(I_0(ta)=tI_0(a)\) for \(0\leq t\leq1\).
If both \(a,ta\) have range in \(K\) and \(t>1\),
apply that identity to \(ta\) with factor \(1/t\).

For arbitrary \(a\in B_0\), take \(M>0\) large
enough that \(a/M\) has range in \([-1,1]\)
and set \(I(a)=M I_0(a/M)\).
Two choices of \(M\) are compared by the just
proved scaling identity, so the answer is
independent of the choice and agrees with
\(I_0\) wherever \(I_0\) was defined.
Common scaling proves monotonicity and
positive homogeneity on all of \(B_0\).
The same construction gives \(I(c\mathbf1)=c\)
for every real constant.

For cash translation, choose a common \(M\)
so that \(2a/M\), \(2c/M\), and all the
quantities used below lie in the central
interval of \(K\). The constant-mixture identity
and scaling give
\[
 I_0((a+c\mathbf1)/M)
       =\tfrac12 I_0(2a/M)+\tfrac12(2c/M)
       =I_0(a/M)+c/M.
\]
Multiplication by \(M\) proves
\(I(a+c\mathbf1)=I(a)+c\).

If \(I(a)=I(b)\), scale both into \(K\).
Their representing acts are indifferent, so
A.5 and positive homogeneity imply
\(I(a+b)\geq I(a)+I(b)\).
For general \(a,b\), put
\(c=I(a)-I(b)\) and \(b'=b+c\mathbf1\).
Cash translation gives \(I(b')=I(a)\).
The equal-value case gives
\(I(a+b')\geq2I(a)\); translating back yields
\(I(a+b)\geq I(a)+I(b)\).
Every arbitrary real simple function can be
scaled into the original range, so these
values are forced by \(J\) and positive
homogeneity. This proves the stated uniqueness
of the extension construction.
\end{proof}
\end{annexitem}

\begin{annexitem}{annex:p17.mmeu.uniform.extension}{thm:p17-mmeu-uniform-extension}{Uniform completion with the algebra unchanged}
\begin{proof}
Bounded functions with the supremum norm
are complete by \cref{thm:p17-sequential-contraction}.
Their closed subspace \(B\) is therefore complete.
If \(d=\|a-b\|_\infty\), then \(a\leq b+d\mathbf1\).
Monotonicity and cash translation give
\(I(a)\leq I(b)+d\), and the reversed inequality
gives the Lipschitz bound.
For any uniformly approximating simple sequence,
these estimates make its \(I\)-values Cauchy.
The same estimate for two approximations proves
independence of the sequence, and defines the
unique continuous extension to \(B\).
Positive homogeneity, superadditivity and
cash translation pass to limits of the
corresponding simple expressions.
If \(a\leq b\), approximations \(a_n,b_n\)
satisfy
\(a_n\leq b_n+(\|a_n-a\|+\|b_n-b\|)\mathbf1\);
the estimate and limit preserve monotonicity.

Products of simple functions are simple after
refining their finite partitions.
Uniformly convergent bounded approximations
satisfy
\(\|a_nb_n-ab\|\leq\|a_n\|\|b_n-b\|+\|b\|\|a_n-a\|\),
so \(ab\in B\).
For continuous \(\phi\) on a compact interval
containing the ranges, uniform continuity gives
\(\phi(a_n)\to\phi(a)\) uniformly.
Each \(\phi(a_n)\) is simple. If necessary
truncate \(a_n\) into that interval first;
truncation preserves simplicity and convergence.
Thus \(\phi(a)\in B\).
All approximations use cells in the original
algebra. No countable union closure of
\(\Sigma\) has been assumed.
\end{proof}
\end{annexitem}

\begin{annexitem}{annex:p17.mmeu.hahn.banach}{thm:p17-mmeu-hahn-banach}{The full dominated extension and normalized support}
\begin{proof}
Suppose \(\ell\) is defined on a subspace \(M\)
and \(\ell\leq p\). For \(z\notin M\), the
interval
\[
 \sup_{x\in M}\{\ell(x)-p(x-z)\}
        \ \leq c\leq\
 \inf_{y\in M}\{p(y+z)-\ell(y)\}
\]
is nonempty. Indeed
\(\ell(x+y)\leq p(x+y)\leq p(x-z)+p(y+z)\)
proves every lower quantity at most every
upper quantity. The choice \(x=y=0\)
bounds the endpoints by finite real numbers,
so a real \(c\) exists.
Set \(\ell'(x+tz)=\ell(x)+tc\).
The decomposition is unique because \(z\notin M\).
For \(t>0\), use the upper inequality at
\(x/t\) and positive homogeneity of \(p\)
to obtain \(\ell'(x+tz)\leq p(x+tz)\).
For \(t<0\), use the lower inequality at
\(x/(-t)\); \(t=0\) is the old domination.
Thus this is a proper one-dimensional extension.

Order all dominated extensions of the initial
functional by restriction. A chain has its
union as a well-defined linear dominated
extension on the union subspace.
The set of such extensions is nonempty.
Zorn's lemma in \cref{thm:p06-ac-to-zorn}
gives a maximal one, and the preceding
one-dimensional extension shows its domain
must be all of \(E\).

Apply this to \(p=-I\), which is sublinear.
On the span of \(b\ne0\), prescribe
\(\ell(tb)=-tI(b)\).
For \(t\geq0\) it equals \(p(tb)\).
For \(t<0\), domination follows from
\(I(b)+I(-b)\leq I(0)=0\).
For \(b=0\) start with the zero subspace.
Extend \(\ell\) and put \(L_b=-\ell\).
Then \(L_b\geq I\) and \(L_b(b)=I(b)\).
If \(a\geq0\), \(I(a)\geq0\), hence \(L_b(a)\geq0\).
The two constants give
\(L_b(\mathbf1)\geq1\) and
\(-L_b(\mathbf1)=L_b(-\mathbf1)\geq-1\);
thus \(L_b(\mathbf1)=1\).
Positivity and
\(-\|a\|\mathbf1\leq a\leq\|a\|\mathbf1\)
give \(|L_b(a)|\leq\|a\|\).
Equality on \(\mathbf1\) makes its norm exactly one.
\end{proof}
\end{annexitem}

\begin{annexitem}{annex:p17.mmeu.charges.compactness}{thm:p17-mmeu-charges-compactness}{Finite additive integrals and the compact prior space}
\begin{proof}
From a positive normalized \(L\), define
\(P(E)=L(\mathbf1_E)\).
Nonnegativity, \(P(S)=1\), and additivity
on disjoint pairs follow from linearity
of indicators. Thus \(P\) is a finitely
additive probability.
For a simple function on a finite disjoint
partition, define its integral by
\(\sum_i a_iP(E_i)\).
Two representations have a common finite
partition refinement, so finite additivity
proves that the sum is representation independent.
It is linear, positive and bounded by
the supremum norm. Uniform approximation
therefore defines a unique integral on \(B\),
with the same properties. Continuity of \(L\)
identifies it with this integral.
Conversely this construction for any \(P\)
gives a positive normalized linear \(L\);
its values on indicators prove uniqueness.
A point mass at any \(s\in S\) shows
the prior space is nonempty.

We prove the compactness premise.
A proper filter on a set is a family
containing the whole set, excluding the
empty set, closed under finite intersections
and supersets. The union of a chain of
proper filters extending a given filter
is again proper. Zorn gives a maximal
one, an ultrafilter \(\mathcal U\).
For every subset \(A\), it contains
\(A\) or its complement: otherwise
adjoining \(A\) preserves the finite
intersection property and produces a
strictly larger proper filter.

Every function into \([0,1]\) has an
ultrafilter limit. Repeatedly choose
a closed half interval whose inverse
image belongs to the ultrafilter,
starting with \([0,1]\). At least
one half works because the two
halves cover the previous interval.
Choose the left half if both work.
Their nested lengths tend to zero;
real completeness gives a unique
point in their intersection.
Every neighborhood of that point
contains a sufficiently small selected
interval, proving the limit.
Disjoint neighborhoods prove uniqueness.

Apply this to every coordinate \(P(E)\)
of an ultrafilter on the prior space.
Let \(P_\infty(E)\) be its scalar limit.
The limits preserve \(P(S)=1,P(\varnothing)=0\)
and finite additivity: intersect the
ultrafilter neighborhoods for \(A,B,A\cup B\)
and let their three error tolerances
tend to zero in the identity
\(P(A\cup B)=P(A)+P(B)\).
Thus \(P_\infty\) is again a prior.
Every basic eventwise neighborhood
constrains finitely many coordinates,
so belongs to the ultrafilter.
For any family of closed subsets
with the finite intersection property,
extend its generated filter to an
ultrafilter. Its limit lies in every
closed member, since otherwise that
member and an open neighborhood of
the limit would be disjoint members
of the filter. This proves compactness:
complements convert failure of a
finite subcover into precisely such
a family. Distinct priors differ
on an event and have disjoint
coordinate neighborhoods, proving
the Hausdorff property.
The same argument applies to a
closed prior subset and to finite
products of compact prior subsets.

For \(a\in B_0\), its integral is
a finite linear combination of
event coordinates and is continuous.
For \(a\in B\), uniform approximation
gives a uniform approximation of
these maps on all priors, since
\(|\int(a-a_n)dP|\leq\|a-a_n\|\).
An epsilon estimate proves continuity
of their uniform limit.
Indicators give the reverse inclusion
of the two generated topologies,
proving the asserted weak-star
identification on the prior space.
For nonempty closed \(C\), let
\(m=\inf_{P\in C}\int a\,dP\).
This is finite, and the closed
nonempty sets
\(\{P\in C:\int a\,dP\leq m+1/n\}\)
have the finite intersection property.
Their intersection supplies a
minimizing prior. No countable
additivity was used.
\end{proof}
\end{annexitem}

\begin{annexitem}{annex:p17.mmeu.support}{thm:p17-mmeu-support}{The exact D1/D2 separator and the prior hull}
\begin{proof}
Superadditivity and positive homogeneity
make \(I\) concave, so \(D_1\) is
convex. Its continuity makes it open,
and a small norm ball about
\(2\mathbf1\) lies inside it.
The two lower sets defining \(D_2\)
are convex; every member of their
convex hull has form
\(\alpha a_1+(1-\alpha)a_2\),
where \(a_1\leq\mathbf1\),
\(a_2\leq b/I(b)\), \(0\leq\alpha\leq1\).
Monotonicity, homogeneity and cash
translation imply
\[
 I(\alpha a_1+(1-\alpha)a_2)
       \leq \alpha+(1-\alpha)I(a_2)\leq1.
\]
Thus \(D_1,D_2\) are disjoint.
The unit ball lies in the first
lower set, so \(D_2\) also has
an interior point.

The supporting functional from the
dominated-extension theorem has
\(L_b\geq I\), \(L_b(\mathbf1)=1\),
\(L_b(b/I(b))=1\), and is positive.
It is consequently greater than
one on \(D_1\) and at most one
on both lower sets and their
convex hull. This is a separator
with exactly the source's normalization
and touching condition; its existence
has been constructed rather than
left as an unnamed separation premise.
The charge correspondence supplies \(P_b\).

Every selected \(P_b\) dominates \(I\)
on all \(B\). Each dominance inequality
is closed and affine in the prior,
so it remains true on their closed
convex hull \(C\), which is compact.
For arbitrary \(a\in B\), choose
a real \(c\) with \(I(a)+c>0\).
The selected prior \(P_{a+c\mathbf1}\)
belongs to the generating set and
satisfies
\[
 \int(a+c\mathbf1)\,dP_{a+c\mathbf1}
                         =I(a+c\mathbf1)=I(a)+c.
\]
Normalization cancels \(c\), so it
attains \(I(a)\). This proves the
representation for all signs of \(I(a)\).
The intrinsic dominance set is also
nonempty, closed and convex, contains
these touching priors, and hence has
the same attained minimum.
\end{proof}
\end{annexitem}

\begin{annexitem}{annex:p17.mmeu.finite.separation}{thm:p17-mmeu-finite-separation}{Finite-coordinate separation on an arbitrary algebra}
\begin{proof}
Because \(P_0\) is outside the closed
set \(C\), a basic eventwise neighborhood
of \(P_0\), using finitely many events
\(E_1,\ldots,E_m\), avoids \(C\).
Project onto these coordinates.
The vector \(x=(P_0(E_i))_i\)
is outside the closed convex set
\(D=\overline{\{(P(E_i))_i:P\in C\}}\);
indeed the neighborhood supplies a
positive distance from this projection.
The set \(D\) is bounded in \([0,1]^m\).
A minimizing sequence for distance to
\(x\) has a convergent subsequence
by bounded-coordinate extraction, so
has a closest point \(y\in D\).
For each \(z\in D\), the segment
\(y+t(z-y)\) lies in \(D\).
Expanding its squared distance to
\(x\) at \(t=0\) gives
\((x-y)\cdot(z-y)\leq0\).
Thus the functional with coefficient
\(v=y-x\) satisfies
\[
               v\cdot x<v\cdot y\leq v\cdot z
                                      \quad(z\in D),
\]
the strict gap being \(\|x-y\|_2^2\).
Set \(a=\sum_i v_i\mathbf1_{E_i}\).
It is simple after refining these
finitely many events, and its
integrals give exactly the displayed
strict separation. This argument also
proves the finite-dimensional
point/closed-convex-set separation
used below.
\end{proof}
\end{annexitem}

\begin{annexitem}{annex:p17.mmeu.representation}{thm:p17-mmeu-representation}{Theorem 1 with both directions and the degenerate exception}
\begin{proof}
Under A.1--A.6, the utility,
certainty-equivalent, simple-functional,
uniform-extension and support results
construct \(u,I,C\), with
\[
       J(f)=I(u\circ f)=\min_{P\in C}\int u\circ f\,dP.
\]
Since \(J\) represents the preference,
this is the forward representation.
If A.6 fails, A.1 gives universal
indifference. A constant utility and
any nonempty closed convex prior set
represent it.

Conversely, every stated minimum is
attained by compactness, since each
\(u\circ f\) is a bounded simple function.
Comparison of the real minimum values
is complete and transitive.
The minimum functional is monotone,
positively homogeneous, cash invariant
and concave: each prior expectation
is linear, and taking its minimum
gives each of these inequalities
directly. It is also 1-Lipschitz
in supremum norm, by comparing
the two expectations under each prior.
The value of mixing with a constant
lottery is exactly the same mixture
of the value and its constant utility.
This proves A.2, including its
strict equivalence. Pointwise lottery
preference means pointwise utility
inequality, proving A.4.
Concavity gives A.5 when the two
values are equal.
For A.3, if \(J(f)>J(g)>J(h)\),
let \(D=\|u\circ f-u\circ h\|>0\).
Choose \(1-\alpha>0\) and \(\beta>0\)
small enough that
\((1-\alpha)D<(J(f)-J(g))/2\)
and \(\beta D<(J(g)-J(h))/2\),
with \(\alpha,\beta\in(0,1)\).
The Lipschitz bound gives
\(J(\alpha f+(1-\alpha)h)>J(g)>
J(\beta f+(1-\beta)h)\).
Thus all five axioms hold.

Utility uniqueness follows from the
constant-lottery representation already
proved. Nonconstant utility gives a
strict constant comparison, and constant
utility makes every act indifferent,
proving the A.6 equivalence.
In the nondegenerate case normalize
two purported representations to
the same \(u\).
If \(P_0\in C_1\setminus C_2\),
finite-coordinate separation gives
a simple \(a\) with
\(\min_{C_1}\int a<\min_{C_2}\int a\).
Scale \(a\) by a positive constant
to have range in \([-1,1]\subseteq K\);
choose an act \(f\in L_0\) with
that utility function. Its two
represented values differ.
But its certainty equivalent in
the common preference is some \(y^*\),
so both values must equal \(u(y)\),
a contradiction. Interchanging the
sets proves uniqueness.

The intrinsic dominance set contains
the constructed \(C\). If it contained
a prior outside \(C\), separation
would give a simple \(a\) whose
integral there is less than
\(\min_C\int a=I(a)\), contradicting
dominance. Hence it equals \(C\).
If utility is constant, any nonempty
closed convex set works. On the
trivial algebra there is only the
prior assigning \(S\) mass one.
On a nontrivial algebra choose a
proper nonempty event and a point
on each side. Their two distinct
Dirac priors give different singleton
representing sets. This proves the
exact degenerate uniqueness boundary.
\end{proof}
\end{annexitem}

\begin{annexitem}{annex:p17.mmeu.extension}{thm:p17-mmeu-extension}{The bounded extension with the original acts retained}
\begin{proof}
Finite-step acts have their level
sets in \(\Sigma\) and constant
upper and lower bounds, so lie
in \(L_{\rm lev}\).
For a level-measurable bounded act,
write \(a=u\circ f\) and bound
it between attained utility values
\(r\leq a\leq s\), \(r,s\in K\).
Finite equally spaced thresholds
between \(r,s\), each represented
by a lottery in \(Y\), give
\(\Sigma\)-simple quantizations
converging uniformly to \(a\).
Thus \(L_{\rm lev}\subseteq L_u\).
For two members of \(L_u\), their
mixed utility is the corresponding
linear combination in \(B\);
the smaller lower bound and
larger upper bound of the two
are again attained in \(K\).
This proves convexity.
Positive affine changes of \(u\)
preserve both properties defining
the domain.

Define the extension by \(I(u\circ f)\).
The same monotonicity, cash, concavity
and Lipschitz proof used in the
representation theorem gives A.1--A.5.
It extends the old preference
because \(I\) agrees with \(J\).
For uniqueness, any competing
axiomatic extension gives every
bounded act a certainty equivalent:
the constant-bounds calibration
uses A.3 exactly as before.
Call its calibrated value \(J'(f)\).
Choose simple \(a_n\to a=u\circ f\),
with error at most \(\epsilon_n\to0\).
Clip \(a_n-\epsilon_n\) and
\(a_n+\epsilon_n\) into \([r,s]\).
The resulting simple functions
\(b_n^-\leq a\leq b_n^+\)
have values in \(K\) and converge
uniformly to \(a\).
Represent them by acts in \(L_0\).
A.4 then bounds \(J'(f)\)
between \(I(b_n^-)\) and
\(I(b_n^+)\). Both limits are
\(I(a)\), proving uniqueness.
The same argument on \(L\cap L_u\)
proves the source's additional
restriction assertion.

If \(\Sigma\) is a sigma algebra,
the uniform limit of simple measurable
functions is measurable: for example
\[
 \{a>c\}=\bigcup_{m\geq1}\bigcup_{N\geq1}
                   \bigcap_{n\geq N}\{a_n>c+1/m\}.
\]
The analogous weak level set follows
by complementation. Thus \(L_u\)
is then \(L_{\rm lev}\).
For a mere algebra this inference
is unavailable, as the next
counterexample demonstrates.
In the degenerate case all
constant lotteries are indifferent;
A.4 in both directions forces
all acts in \(Y^S\) indifferent,
giving the stated extension.
\end{proof}
\end{annexitem}

\begin{annexitem}{annex:p17.mmeu.domain.boundaries}{thm:p17-mmeu-domain-boundaries}{All algebra and extension counterexamples}
\begin{proof}
On the finite/cofinite algebra,
every simple function is eventually
constant. Its uniform completion
is exactly the convergent sequences:
uniform approximation by eventually
constant sequences makes the tail
Cauchy, and a convergent sequence
is approximated by retaining a
finite prefix and replacing its
tail by its limit.
The displayed \(a(n)\) approaches
\(1/2\) from above and \(b(n)\)
from below. Each strict or weak
upper level set is finite or
cofinite; at level \(1/2\) it
is respectively the whole set
or the empty set. Both functions
take values in \([0,1]\), so
are legitimate binary lottery acts.
Their half mixture has utility
\(1/2+(-1)^n/(16n^2)\).
Its upper level set at \(1/2\)
is the infinite coinfinite even
set, proving failure of convexity
of \(L_{\rm lev}\). All three
utility functions converge,
so belong to \(B\) and \(L_u\).

Extend the cofinite filter on
\(\mathbb N_+\) to an ultrafilter
by the proved filter construction.
It contains no finite set.
An ultrafilter contains a disjoint
union if and only if it contains
one of the two disjoint members,
and cannot contain both.
Thus \(P_{\mathcal U}\) is a
finitely additive probability.
All singletons have mass zero,
while their countable union has
mass one. The function \(1/n\)
is uniformly approximated by
finite-support sequences, each
having integral zero, so its
integral is zero.

Define \(f\succeq_{\mathcal U}g\)
when \(\{n:u(f(n))\geq u(g(n))\}\)
belongs to this ultrafilter.
Ultrafilter decisions give completeness;
finite intersection and transitivity
of the pointwise order give
transitivity. Pointwise improvement
gives the whole set, so the order
is monotone. On simple acts,
utilities are eventually constant,
and this comparison is exactly
the comparison of those constants.
But \(1/n\succ_{\mathcal U}0\),
and \(1\succ_{\mathcal U}1/n\).
For every \(\beta>0\), the set
\(\{n:1/n>\beta\}\) is finite,
so \(1/n\) is not preferred
to the constant \(\beta\).
A.3 for \(1\succ1/n\succ0\)
therefore fails. The canonical
integral extension instead makes
\(1/n\) indifferent to zero.
This disproves uniqueness under
monotonicity alone.

Finitely supported priors on
all subsets of \(\mathbb N_+\)
form a convex set. A simple
function attains its smallest
value on some nonempty cell,
so a point mass there attains
its minimum under both this
set and the whole prior space.
For every finite collection
of events, their finite Boolean
partition has one cell belonging
to the ultrafilter. A point in
that cell gives a Dirac prior
agreeing with \(P_{\mathcal U}\)
on all these events. Hence
every basic neighborhood of
\(P_{\mathcal U}\) meets the
finitely supported priors.
That prior is not finitely
supported, so the set is
not closed. This proves the
claimed necessity of the
closedness qualification.
\end{proof}
\end{annexitem}

\begin{annexitem}{annex:p17.mmeu.extreme.geometry}{thm:p17-mmeu-extreme-geometry}{The complete compact-convex geometry}
\begin{proof}
Order the nonempty closed faces of a nonempty
closed face \(F\subseteq K\) by reverse inclusion.
A chain has nonempty intersection by compactness
and the finite intersection property; that intersection
is again a closed face. Zorn supplies a minimal
nonempty closed face \(G\).
If \(P,Q\in G\) differ, some event coordinate
has different values. Its maximum on \(G\)
is attained, and the maximizing set is a
proper nonempty closed face of \(G\), hence
of \(F\) and \(K\), a contradiction.
Thus \(G\) is a singleton and gives an
extreme point of \(K\) in \(F\).
The assertion that a face of a face is a
face follows directly by applying the
defining convex-combination implication twice.

Let \(H=\overline{\operatorname{conv}}\operatorname{Ext}(K)\).
It is nonempty, closed, convex and contained in \(K\).
If \(P_0\in K\setminus H\), finite-coordinate
separation gives a simple \(a\) with
\(\int a\,dP_0<\min_H\int a\,dP\).
The minimum of this functional on \(K\)
is attained on a nonempty closed face.
That face has an extreme point \(Q\) of \(K\),
so
\[
 \int a\,dQ=\min_K\int a\,dP
       \leq\int a\,dP_0<\min_H\int a\,dP
       \leq\int a\,dQ,
\]
a contradiction. This proves the closed-hull assertion.

For a continuous affine map \(\pi\), the
fiber over an extreme point \(y\in\pi(K)\)
is a nonempty closed face of \(K\):
if a strict convex combination maps to \(y\),
extremeness forces both images to be \(y\).
The first result supplies an extreme point
of \(K\) in that fiber, proving the lift.

Finally suppose \(K=\overline{\operatorname{conv}}D\)
with \(D\) compact, and let \(P_0\) be
extreme in \(K\) but outside \(D\).
A compact set in a Hausdorff space is closed:
separate an outside point from each point
of the compact set and take a finite
subcover of the latter neighborhoods.
Thus a finite-coordinate neighborhood of
\(P_0\) avoids \(D\).
There are finitely many coordinates and
an \(\epsilon>0\) such that every member
of \(D\) differs from \(P_0\) in one
of them by at least \(\epsilon\).
Cover \(D\) by the resulting finitely
many closed coordinate halfspace pieces
\(D_1,\ldots,D_m\), omitting empty pieces.
Put \(K_j=\overline{\operatorname{conv}}D_j\).
These are nonempty compact subsets of \(K\),
and none contains \(P_0\), since the
corresponding halfspace inequality persists.

The set of sums \(\sum_j\lambda_jQ_j\),
\(\lambda\in\Delta_m,Q_j\in K_j\), is compact
and hence closed: the finite product is
compact by the same ultrafilter-coordinate
argument used for priors, and the sum map
is continuous. It contains every finite
convex combination of points of \(D\),
by assigning each point to a piece and
grouping its coefficients. It therefore
contains \(K\), in particular \(P_0\).
But an extreme point of \(K\) represented
by such a finite convex combination must
equal every positively weighted \(Q_j\).
For a weight strictly between zero and
one, combine the other terms into one
point of \(K\) and use extremeness;
weight one gives equality directly.
This contradicts \(P_0\notin K_j\).
Thus every extreme point belongs to \(D\).
\end{proof}
\end{annexitem}

\begin{annexitem}{annex:p17.mmeu.exposed.points}{thm:p17-mmeu-exposed-points}{The exposed-point consequence with a direct finite-dimensional proof}
\begin{proof}
A singleton is immediate. Otherwise fix
a vector \(v\) and a positive number \(t\).
Choose \(x_t\in K\) farthest from \(-tv\),
which exists by compactness.
For any \(z\in K\),
\[
 \|z+tv\|_2^2\leq\|x_t+tv\|_2^2
 \quad\Longrightarrow\quad
 (x_t+tv)\cdot(z-x_t)\leq-\tfrac12\|z-x_t\|_2^2.
\]
The normal \(x_t+tv\) is nonzero when
\(K\) has more than one point.
It therefore exposes \(x_t\) uniquely.
If \(y\) maximizes \(v\cdot x\) on \(K\),
the farthest-point inequality also gives
\[
 0\leq v\cdot y-v\cdot x_t
       \leq\frac{\max_{z\in K}\|z\|_2^2
                      -\min_{z\in K}\|z\|_2^2}{2t}.
\]
Thus exposed points approximate every
linear maximum. If a point of \(K\)
were outside the closed convex hull
of exposed points, the finite-dimensional
nearest-point separation proved above
would strictly separate it from that
hull. This contradicts the just proved
equality of the two suprema for that
linear functional. The hull is contained
in compact \(K\), so its closure is
compact and the separating argument
applies. This proves precisely the
exposed-point hull fact needed below,
without leaving an approximation
theorem unproved.
\end{proof}
\end{annexitem}

\begin{annexitem}{annex:p17.mmeu.product.priors}{thm:p17-mmeu-product-priors}{Product charges and both source independence clauses}
\begin{proof}
Finitely many rectangles can be refined
by the finite Boolean partitions generated
by their sides in each marginal algebra.
Assign a cell \(A_i\times B_j\) mass
\(P_1(A_i)P_2(B_j)\) and sum over its
selected cells. Finite additivity in
each coordinate proves invariance under
further refinements, so defines a
well-defined nonnegative finitely additive
probability on the rectangle algebra.
Any such probability is determined by
these rectangle values, proving uniqueness.
For a fixed event the resulting value
is a finite sum of products of marginal
event coordinates, hence is continuous.
The finite product \(C_1\times C_2\)
is compact by the ultrafilter argument;
its continuous image \(D\) of product
priors is compact.
Its closed convex hull \(C\) lies in
the compact prior space and is
nonempty, compact and convex.
The final assertion of the extreme
geometry theorem gives
\(\operatorname{Ext}(C)\subseteq D\).

The marginal map \(R_1(P)(A)=P(A\times S_2)\)
is continuous and affine. The closed
convex set \(R_1^{-1}(C_1)\) contains
all generators, so contains \(C\).
Conversely every \(P_1\in C_1\)
occurs as a marginal of its product
with a fixed member of \(C_2\).
Thus \(R_1(C)=C_1\); the other
marginal is the same argument.
The identity
\(\int\bar a\,d(P_1\otimes P_2)=\int a\,dP_1\)
holds first for simple functions and
then by uniform approximation for
all \(a\in B(\Sigma_1)\).

For two lifted acts, choose a prior
minimizing each marginal utility
integral. Their product belongs to
\(C\) and minimizes both lifted
integrals, by the marginal identities.
This proves clause 1.
Every extreme joint prior is a product.
A measurable event of a lifted
utility function has form \(A\times S_2\)
or \(S_1\times B\). If such a
cylinder belongs to the rectangle
algebra, its section at any point
of the other nonempty state set
belongs to the corresponding marginal
algebra, by finite rectangle refinement.
Thus the product formula applies to
every pair of events specified in
clause 2, proving independence.
Positive affine utility changes preserve
all minimizing priors and merely
relabel Borel utility sets by an
affine bijection. Both clauses are
therefore invariant, as asserted.
\end{proof}
\end{annexitem}

\begin{annexitem}{annex:p17.mmeu.product.representation}{thm:p17-mmeu-product-representation}{The full product uniqueness proof and its finite restrictions}
\begin{proof}
The marginal constant orders agree
and are nondegenerate. Utility uniqueness
therefore permits one common normalized
\(u\). The product-prior construction,
representation theorem and bounded
extension give a preference on \(L_u\)
satisfying A.1--A.6, retaining both
marginal comparisons and both independence
clauses. This proves existence.

Let a competing preference satisfy
the three conditions. Its restriction
to product simple acts is nondegenerate,
since the inherited constant order is.
The representation and extension
theorems give it the same normalized
utility and a unique nonempty closed
convex joint prior set \(C'\).
Its marginal images are compact,
closed and convex. They represent
the prescribed marginal preferences
on simple acts, so marginal prior
uniqueness gives exactly \(C_1,C_2\).

For \(Q\in\operatorname{Ext}(C')\),
take binary marginal acts with
utilities \(\mathbf1_A,\mathbf1_B\)
for arbitrary \(A\in\Sigma_1,B\in\Sigma_2\).
The normalization of \(u\) permits
those utility values.
Clause 2 gives
\[
                Q(A\times B)
                   =Q(A\times S_2)Q(S_1\times B).
\]
Thus \(Q\) is the product of its
marginals, which belong to \(C_i\).
The extreme-point closed-hull theorem
therefore implies \(C'\subseteq C\).

First suppose the two algebras are
finite. Regard their priors as
vectors of masses on their nonempty
atoms. Let \(P_i\) be exposed
points of \(C_i\). Choose linear
functionals uniquely minimized at
these points, reversing an exposing
sign if necessary. Scale their
finitely many coefficients into
the central interval of \(K\)
and realize them by simple
marginal acts \(f,g\).
By clause 1, the two minimizing
faces of \(C'\) for the lifted
integrals have nonempty intersection.
That intersection is a closed
face of \(C'\): equality in a
minimum of a linear functional
forces equality at both endpoints
of a strict convex combination.
Choose an extreme point \(Q\)
of \(C'\) in that common face,
using the proved face theorem.
The preceding paragraph makes
this particular \(Q\) a product.
Its marginals uniquely minimize
the two exposing functionals,
so are \(P_1,P_2\). Hence
\(P_1\otimes P_2\in C'\).

Each \(C_i\) is the closed convex
hull of its exposed points by
the finite-dimensional theorem.
Approximate arbitrary \(P_i\in C_i\)
by finite convex combinations of
exposed points. Their products
are convex combinations of the
exposed products just shown to
belong to \(C'\), by distributing
the two finite sums. Continuity
of the finite-dimensional product
and closedness of \(C'\) give
\(P_1\otimes P_2\in C'\).
Taking the closed convex hull
proves \(C\subseteq C'\) in the
finite-algebra case.
The use of an extreme point in
the common minimizing face is
essential: an arbitrary point of
\(C'\subseteq C\) need not itself
be a product.

Now take arbitrary algebras.
Choose finite subalgebras
\(\Sigma_i^0\subseteq\Sigma_i\)
and restrict joint priors to
\(\Sigma_1^0\otimes\Sigma_2^0\).
The image \(R\) of \(C'\) is
compact and convex, and its
marginal images are the corresponding
images \(C_i^0\) of \(C_i\).
Clause 1 passes to these restrictions
because the finite-subalgebra acts
lift to original acts with the
same integrals.
Every extreme point of \(R\)
lifts to an extreme point of
\(C'\), by the proved extreme
projection theorem; the latter
is a product. Thus clause 2
holds for the restricted priors.
The finite-algebra proof gives
\[
             R=\overline{\operatorname{conv}}
                       \{P_1^0\otimes P_2^0:P_i^0\in C_i^0\}.
\]
This is also the image of \(C\).
Indeed that image is closed
convex and contains all these
products; conversely the inverse
image of their closed convex
hull is a closed convex set
containing all full product
generators and therefore \(C\).

If \(C\not\subseteq C'\), separate
a point of \(C\setminus C'\)
from \(C'\) by finitely many
event coordinates as above.
Each of these product-algebra
events is a finite union of
rectangles. The finitely many
marginal sides generate finite
subalgebras \(\Sigma_i^0\).
Equality of the two restricted
images supplies a member of \(C'\)
with exactly the same selected
coordinates, contradicting that
separation. Hence \(C=C'\)
on the original arbitrary algebras.
The representation and extension
theorems then identify the two
preferences on all of \(L_u\).
This proves the entire characterization.
\end{proof}
\end{annexitem}

\begin{annexitem}{annex:p17.mmeu.product.examples}{thm:p17-mmeu-product-examples}{The exact product values and two separate failure modes}
\begin{proof}
A binary product with success
probabilities \(p,q\) has vector
\(((1-p)(1-q),(1-p)q,p(1-q),pq)\).
At the four endpoint pairs
\(p,q\in\{1/4,3/4\}\), these
are the four listed vectors.
Writing each marginal prior as
a convex combination of its
two endpoints and distributing
their product shows that every
product in the two intervals
lies in this four-point convex
hull. Each endpoint product
is extreme: a decomposition
in the hull must preserve an
extreme endpoint of each marginal
interval, forcing both marginals
and hence the corner itself.

Half the first and last vectors
is \((5,3,3,5)/16\).
Its marginals are both fair,
but its mass at \(11\) is
\(5/16\ne1/4\), so it is
not a product. Linear minimization
over the four corners gives
each marginal success value
\(1/4\) and the joint success
value \(1/16\).

The hull of the middle two
products has both marginal
images \([1/4,3/4]\), and
its extreme points are products.
The first success probability
is minimized only at the
corner with \((p,q)=(1/4,3/4)\);
the second only at the other
corner. On their segment
\(p+q=1\), so they cannot
both equal \(1/4\). There is
no common minimizing prior.

For singleton fair marginals,
take the segment joining
\((1/2,0,0,1/2)\) and
\((0,1/2,1/2,0)\).
Every lifted marginal act
has the same expectation at
every member, so every prior
is a simultaneous minimizer.
Both endpoints are extreme,
but their \(11\) probabilities
are \(1/2\) and zero, neither
equal to the product value
\(1/4\). Thus the extreme-point
independence condition fails.
These are represented preferences
on the full finite lottery-act
domains, satisfying A.1--A.6,
so the examples isolate the two
additional product conditions.
\end{proof}
\end{annexitem}
